%% ============================================================
%%  A2 — Phase 3 Signal Research: BTCUSDT Perpetual Futures
%%  Compile with: pdflatex A2.tex  (twice for refs)
%% ============================================================
\documentclass[12pt, a4paper]{article}

\usepackage{amsmath, amssymb}
\usepackage{booktabs}
\usepackage{hyperref}
\usepackage{geometry}
\usepackage{enumitem}

\geometry{margin=2.5cm}
\hypersetup{colorlinks=true, linkcolor=blue!60!black, urlcolor=blue!60!black}

\newcommand{\IC}{\mathrm{IC}}
\newcommand{\DVOL}{\mathrm{DVOL}}
\newcommand{\OI}{\mathrm{OI}}
\newcommand{\VRP}{\mathrm{VRP}}
\newcommand{\TAR}{\mathrm{TAR}}
\newcommand{\bp}{\,\mathrm{bp}}

\title{\textbf{Phase 3 Signal Research --- BTCUSDT Perpetual Futures}\\[0.3em]
\large Volatility Risk Premium, Taker Aggression, and OI-Price Divergence}
\author{}
\date{May 2026}

\begin{document}
\maketitle

\begin{abstract}
We test three candidate daily-frequency signals against BTCUSDT perpetual
futures returns using the same 7-stage validation pipeline applied in report
A1. Signal P3 (OI-Price Divergence, DD regime) passes all validation gates:
OOS Spearman IC $= +0.054$ ($p_\mathrm{boot} = 0.032$, $n=819$ days),
position-level OOS Sharpe $+5.37$, and block-bootstrap $p = 0.007$ on
57 exclusive trades where N3 was simultaneously inactive. Signal correlation
with N3 is $+0.027$, establishing near-zero overlap. P3 is deployed to
shadow paper trading. Signal P1 (Volatility Risk Premium) shows genuine IC
at 48--72h but collapses when N3 activity is excluded (P1-exclusive Sharpe
$= -0.50$), indicating P1 proxies N3's fear-regime driver rather than
contributing independent information. Signal P2 (Taker Aggression Ratio) has
no detectable daily-frequency signal. All results are conditional on the
stated cost model and data; specific entry thresholds are withheld from this
public version.
\end{abstract}

%% ============================================================
\section{Introduction}
%% ============================================================

Report A1 validated signal N3 --- the 30-day DVOL z-score with a volatility
regime filter --- achieving OOS Sharpe $+2.95$ over 863 days (2024--2026
YTD). Phase 3 tests three structurally distinct signal families.
\textbf{P1} hypothesises that implied volatility substantially exceeding
recent realised volatility predicts positive returns via mean-reversion of the
fear premium. \textbf{P2} tests whether aggressive buyer-initiated order flow
at daily frequency predicts returns; the prior from market microstructure
theory is negative, as informed flow aggregates away at this timescale.
\textbf{P3} hypothesises that simultaneous price decline and OI contraction
identifies long-liquidation exhaustion: forced sellers are bounded, and the
supply vacuum after liquidation cascades clear produces a predictable
mean-reversion impulse.

%% ============================================================
\section{Data and Methodology}
%% ============================================================

All signals are evaluated at daily frequency (non-overlapping 00:00 UTC
samples) over 2023--2026-YTD using Binance FAPI 1m OHLCV, 5m open interest,
8h funding rates, and Deribit BTCVOL 1h. Data splits, the Spearman IC
criterion, breakeven IC derivation ($\IC^* \approx 0.035$ under 6 bp
round-trip maker cost), overlapping-return adjustment, and block bootstrap
(21-day blocks, $B=2{,}000$, one-sided) are documented in full in A1
Sections 3--5. No signal parameters are fitted using OOS data. Train: 2023;
OOS: 2024-H1 through 2026-YTD.

%% ============================================================
\section{Results}
%% ============================================================

\subsection{P1: Volatility Risk Premium}

VRP is defined as $\DVOL_t - \sigma^{\mathrm{rv}}_t$ (30-day realised vol),
z-scored over a 30-day rolling window. The signal is statistically
significant at 48--72h in the full sample ($p_\mathrm{boot} < 0.025$, IC
$= +0.075$--$+0.086$). However, it is strongly regime-conditional: IC
$= +0.140$ in the DVOL $51$--$60$ band ($p_\mathrm{boot} < 0.0005$) but
near zero outside it. The deep dive in Section~\ref{sec:p1indep} shows this
conditional significance is entirely attributable to N3 co-firing days.
\textbf{P1 is not promoted.}

\subsection{P2: Taker Aggression Ratio}

The daily taker-buy fraction $\TAR_t = V^{\mathrm{buy}}_t / V_t$ is near
$0.50$ on almost every day ($\mu = 0.497$, $\sigma = 0.013$), confirming
that informed order flow aggregates away at daily resolution. Across twelve
signal-horizon combinations (TAR, TAR$^z$, $\Delta_5$TAR at 8h, 24h, 48h,
72h), only one isolated result reaches $p < 0.05$ and does not survive any
multiple-testing correction. \textbf{P2 is killed at the IC screen.}

\subsection{P3: OI-Price Divergence (DD Regime)}

The DD regime ($\Delta P < 0$, $\Delta\OI < 0$) identifies long-liquidation
exhaustion. The four regimes are approximately uniformly distributed; the DD
subset is not selected by unusual frequency. The IC screen confirms the
signal's predictive content:

\begin{center}
\begin{tabular}{lrrrr}
\toprule
Horizon & $n$ & IC & $p$-boot & Sig. \\
\midrule
24h & 1183 & $+0.054$ & 0.032 & $*$ \\
48h & 1183 & $+0.059$ & 0.023 & $*$ \\
72h & 1183 & $+0.059$ & 0.019 & $*$ \\
\midrule
OOS confirmed & 819 & $+0.065$ & 0.026 & $*$ \\
\bottomrule
\end{tabular}
\end{center}

The position-level backtest (24h hold, maker cost, DVOL regime filter)
shows consistent results across the full OOS period:

\begin{center}
\begin{tabular}{lrrrrrr}
\toprule
Period & $n$ & Sharpe & PnL (bp) & MaxDD (bp) & Win\,\% & $p$-boot \\
\midrule
Train 2023   & 34 & $+5.49$ & $+2{,}532$ & $-643$ & 58.8 & 0.006 \\
OOS 2024-H1  & 34 & $+5.51$ & $+3{,}691$ & $-658$ & 67.6 & $<0.001$ \\
OOS 2024-H2  & 35 & $+7.52$ & $+3{,}232$ & $-530$ & 68.6 & $<0.001$ \\
OOS 2025-H1  & 19 & $+3.11$ & $+1{,}029$ & $-642$ & 68.4 & 0.295 \\
OOS 2025-H2  & \multicolumn{6}{l}{\textit{$<$2 trades --- DVOL below $\delta$ throughout H2 2025 bull market}} \\
OOS 2026-YTD & 10 & $+5.57$ & $+1{,}616$ & $-518$ & 50.0 & 0.397 \\
\bottomrule
\end{tabular}
\end{center}

Sharpe is positive in all five periods with sufficient trade count.
Removing the two most extreme return days or all known macro-event days
changes the overall Sharpe by less than 0.6 (99 $\to$ 97 trades,
Sharpe $+5.37 \to +4.83$), indicating the edge is distributed across the
calendar rather than concentrated in tail events. The signal improves
monotonically with DVOL: IC $= +0.005$ at DVOL $< 51$,
$+0.051$ at $51$--$60$, and $+0.180$ at DVOL $\geq 60$
($p_\mathrm{boot} < 0.0005$), consistent with liquidation cascades being
most severe in high-fear environments.

\subsection{P3 Independence from N3}
\label{sec:p3indep}

The independence question is whether P3 adds signal content beyond what N3
already captures. Spearman correlation between the two binary signals is
$+0.027$ (statistically indistinguishable from zero); they fire on almost
entirely different days (only 42 of 226 DD days overlap with an N3 long
signal). The key evidence is the performance on \emph{exclusive} P3 trades:

\begin{center}
\begin{tabular}{lrrrrrr}
\toprule
Strategy & $n$ & Sharpe & PnL (bp) & MaxDD (bp) & Win\,\% & $p$-boot \\
\midrule
N3 long-only              & 146 & $+3.38$ & $+10{,}094$ & $-1{,}150$ & 54.1 & $<0.001$ \\
P3 (DD) only              &  99 & $+5.37$ & $+9{,}510$  & $-972$     & 65.7 & $<0.001$ \\
\textbf{P3 excl.\ N3}     & \textbf{57} & $\mathbf{+5.18}$ & $+4{,}599$ & $-769$ & 68.4 & \textbf{0.007} \\
P3 OR N3                  & 203 & $+3.79$ & $+14{,}741$ & $-1{,}293$ & 58.1 & $<0.001$ \\
\bottomrule
\end{tabular}
\end{center}

On the 57 trades where N3 was completely inactive, P3 achieves Sharpe $+5.18$
with $p_\mathrm{boot} = 0.007$. This is the primary statistical evidence that
P3 carries independent signal content. The combined ``P3 OR N3'' strategy
produces $+14{,}741\bp$ OOS cumulative PnL, the highest of any configuration,
confirming complementarity rather than duplication.

\subsection{P1 Independence from N3}
\label{sec:p1indep}

The critical test: when N3 is excluded from the DVOL $51$--$60$ band
(the only region where P1 is significant), the VRP-exclusive Sharpe is
$-0.50$ ($p_\mathrm{boot} = 0.732$, $n = 75$ trades). The VRP edge in
this band is entirely driven by days where N3z is also elevated; both signals
respond to the same fear-regime driver. Additionally, OOS 2025-H1 shows
Sharpe $-0.41$ ($n = 23$) despite sufficient trade count, a period where
P3 remained positive ($+3.11$). \textbf{P1 does not add independent
information when N3 is inactive.}

%% ============================================================
\section{Summary and Deployment}
%% ============================================================

\begin{center}
\begin{tabular}{llrlrl}
\toprule
Signal & Screen IC & $p$-boot & Backtest Sh & Indep.\ from N3 & Decision \\
\midrule
P1 VRP   & $+0.075$ (48h) & 0.017 & $+3.65$ & No ($p=0.73$)   & Not promoted \\
P2 TAR   & $+0.013$ (24h) & 0.304 & ---     & ---             & \textbf{Killed} \\
P3 OI-PD & $+0.054$ (24h) & 0.032 & $+5.37$ & Yes ($p=0.007$) & \textbf{Deployed (shadow)} \\
\midrule
N3 (A1)  & $+0.112$ (24h) & 0.039 & $+3.38$ & ---             & Validated \\
\bottomrule
\end{tabular}
\end{center}

P3 was promoted to shadow paper trading on 14 May 2026 under the frozen rule:
\[
  \Delta p_{24h} < 0 \;\wedge\; \Delta\OI_{24h} < 0 \;\wedge\; \DVOL \geq \delta
  \quad\Longrightarrow\quad \text{LONG},\; 24\mathrm{h},\; \text{maker}
\]
\noindent\textit{(Threshold $\delta$ is withheld from this public version.)}

Forward review protocol ($\geq$ 14 Aug 2026): promote if Sharpe $> 0$ on
$\geq$15 exclusive P3 trades; reassess if Sharpe $< 0$ on $\geq$20 exclusive
trades. The ``P3 excl.\ N3'' row in the performance API is the primary monitor.

%% ============================================================
\section{Conclusion}
%% ============================================================

P3 is a statistically independent second candidate signal from N3. The
exhaustion-bounce mechanism --- forced long liquidations creating a supply
vacuum --- produces robust OOS results with win rates of 65--72\% and
positive Sharpe in four of five measurable OOS periods. Signal correlation
with N3 of $+0.027$ and exclusive-trade $p = 0.007$ are the primary
independence evidence. P1 fails the independence test despite genuine IC;
P2 has no daily-frequency signal. The DU short-side hypothesis was also
tested and killed (bootstrap $p = 0.474$; negative Sharpe in 8 of 9
grid cells, 5 of 6 half-year periods negative).

\textit{A conservative reading: P3's Sharpe of $+5.18$ rests on 57 exclusive
trades. Crypto Sharpe figures are sensitive to the cost model and OOS window;
these results should be treated as a strong candidate signal under the stated
assumptions, not a confirmed persistent edge. The 3-month forward review will
provide the first genuinely prospective evidence.}

\vspace{0.5em}
\noindent\textit{Data: Binance FAPI 1m klines, 5m OI, 8h funding (2022--2026);
Deribit BTCVOL 1h (2022--2026). Entry thresholds withheld from public version.}

%% ============================================================
\section*{References}
%% ============================================================
\begingroup
\setlength{\parindent}{-2em}
\setlength{\leftskip}{2em}
\setlength{\parskip}{4pt}
\small

Almeida, C., Grith, M., Miftachov, R., \& Wang, Z. (2024). Risk premia in the Bitcoin market. \textit{arXiv preprint} arXiv:2410.15195.

Bollerslev, T., Tauchen, G., \& Zhou, H. (2009). Expected stock returns and variance risk premia. \textit{The Review of Financial Studies}, \textit{22}(11), 4463--4492.

Cont, R., Kukanov, A., \& Stoikov, S. (2014). The price impact of order book events. \textit{Journal of Financial Econometrics}, \textit{12}(1), 47--88.

Glosten, L. R., \& Milgrom, P. R. (1985). Bid, ask and transaction prices in a specialist market with heterogeneously informed traders. \textit{Journal of Financial Economics}, \textit{14}(1), 71--100.

Künsch, H. R. (1989). The jackknife and the bootstrap for general stationary observations. \textit{The Annals of Statistics}, \textit{17}(3), 1217--1241.

Kyle, A. S. (1985). Continuous auctions and insider trading. \textit{Econometrica}, \textit{53}(6), 1315--1335.

\endgroup

\end{document}
